\documentclass[11pt,a4paper]{article}
\usepackage[T1]{fontenc}
\usepackage[utf8]{inputenc}
\usepackage[brazil]{babel}
\usepackage{lmodern}
\usepackage{microtype}
\usepackage[a4paper,margin=2.5cm]{geometry}
\usepackage{parskip}
\pagestyle{empty}

\begin{document}
\begin{center}
{\large Projeto Talos: cofre eletrônico com ESP32}\\[4pt]
PCS3848 Sistemas Embarcados, Escola Politécnica da USP\\[4pt]
Autoavaliação de André Moreno (13682785)
\end{center}
\vspace{0.6cm}

O grupo cumpriu o essencial do planejamento. O cofre já executa o fluxo completo na placa, com desafio anti-marcas, validação da senha por hash, bloqueio progressivo, troca de senha, modo stealth e deep sleep com despertar por qualquer tecla. A evolução física avançou até o protótipo em bancada com o teclado, o buzzer e a fechadura testada na bateria de 9 V. A montagem no case e o display ainda dependem de peças e de tempo de bancada. A documentação acompanhou o trabalho, com o guia de montagem, o guia de testes e o relatório técnico.

Minha contribuição foi o desenvolvimento de todo o firmware do cofre. Escrevi a máquina de estados, a leitura do teclado, a persistência com cópia dupla e CRC, o acionamento da fechadura e o driver do display. Também conduzi os testes na placa e resolvi os problemas de mapeamento do teclado, de gravação na memória e de despertar pelo pino D12.

Atribuo a mim a nota 10.

\end{document}
